\section{Česte zamke}\label{sec:zamke}

\begin{enumerate}[leftmargin=1.6em]
\item \textbf{$\alpha=b$.} Metabolički eksponent miješa tokove koji skaliraju s $V^{2/3}$, $V$ i $V^{\alpha}$ te ovisi o udjelu rezerve u masi (T1, T2). Isti raspon $[2/3,1]$ za $b$ (T12) i za $\alpha$ posljedica je istog miješanja površinskih i volumnih procesa, a ne dokaz da su jednaki.
\item \textbf{Masa umjesto strukture.} MTE izražava skaliranje prema ukupnoj masi $M$ (T7–T9), a DEB razlikuje strukturu i rezervu i skalira prema strukturi (T1, T2). Ako gustoća spremišta ovisi o veličini (P2), eksponenti po masi i po strukturi se razlikuju. Isto vrijedi za udio lipida: udio po ukupnoj masi $\varphi$ treba pretvoriti u $\varphi/(1-\varphi)$ (pogl.~\ref{sec:protokol}, Korak 3).
\item \textbf{Interspecifično umjesto intraspecifičnog.} Interspecifični $3/4$ uključuje promjenu parametara među vrstama (T2). $\alpha$ je ontogenetski eksponent.
\item \textbf{Kisik i energija kao jedna mreža.} Traheje kukaca i kožno disanje vodozemaca odvajaju kisik od energije (H5).
\item \textbf{Morfometrija kao procjena.} Strukture prate vršnu potražnju i skaliraju strmije od bazalnog metabolizma (E1, E2, T15). Morfometrija daje prior, ne $\alpha$.
\item \textbf{Fraktalna dimenzija sa slika.} Dimenzija manja od dimenzije prostora znači da mreža \emph{ne} ispunjava prostor, pa je $\alpha<1$ (H4, P6).
\item \textbf{Nagib bez dinamike.} Spremište zaostaje za rastom, a naivna procjena pristrana je za \NaiveBiasMin–\NaiveBiasMax{} (S4).
\item \textbf{Jedan $\alpha$ za cijeli život.} Kod riba se mreža izmjene mijenja s kože na škrge (E4, E5).
\item \textbf{Fitanje samo na rast.} $\alpha$ tada nije identifikabilan (P1, S2, S6).
\item \textbf{Pretpostavljeni umjesto izmjerenog $a$.} Ako je istinski $a$ veći od pretpostavljenog $2/3$, $\hat\alpha$ je pristran naniže za približno $a-2/3$ (P2), i to i uz podatke o gladovanju kad su u fitu i spremište i rast (S6).
\item \textbf{$\beta$ bez normalizacije.} Jedinice $\beta$ ovise o $\alpha$, pa se $\beta$ ne mogu uspoređivati među vrstama s različitim $\alpha$ (pogl.~\ref{sec:definicija}). Koristiti jedn.~\eqref{eq:norm}.
\item \textbf{Parametri bez dokumentiranog izvora.} Ilustrativne vrijednosti, poput onih u pogl.~\ref{sec:simulacije}, lako se prenesu u analizu kao da su procjene za stvarnu vrstu. Svaki parametar treba imati dokumentiran izvor: mjerenje, literaturu s navedenim podacima ili fit s navedenim podacima i postupkom.
\end{enumerate}
